\section*{Hoofdstuk \ref{ch:b3:computational-chemistry} --- Computationele chemie: Hartree-Fock en DFT}

\begin{solution}{exo:b3:computational-chemistry:1}
$\dd E/\dd\alpha = \frac32 - \sqrt{2/\pi\alpha} = 0$ geeft $\alpha = 8/9\pi = 0.283\,
a_0^{-2}$ en $E = \frac32\alpha - 2\sqrt{2\alpha/\pi} = -4/3\pi = \qty{-0.4244}{\hartree}$,
15\,\% boven de exacte $-0.5$: één enkele gaussfunctie heeft noch de knik
noch de staart van de $1s$-functie.
\end{solution}

\begin{solution}{exo:b3:computational-chemistry:2}
STO-3G: 5 functies per C ($1s$, $2s$, drie $2p$) en 1 per H: $6 \times 5 + 6 =
36$. 6-31G(d): per C 1 (kern) + $2 \times 4$ (gesplitste valentie) + 6 ($d$) $= 15$;
per H 2: $6 \times 15 + 6 \times 2 = 102$.
\end{solution}

\begin{solution}{exo:b3:computational-chemistry:3}
De elektronische \omterm{def:b3:quantum-model-systems:schrodinger}{hamiltonoperator} bevat geen kernmassa's: binnen de
\omterm{def:b3:computational-chemistry:born-oppenheimer}{born-oppenheimerbenadering} zijn het PES, en dus $r_e$ en $k = U''(r_e)$,
hetzelfde. $\tilde\omega = \sqrt{k/\mu}/2\pi c$ hangt af van de
\omterm{def:b3:quantum-model-systems:oscillator}{gereduceerde massa}, die verdubbelt: verhouding $\sqrt2 = 1.414$ (gemeten 1.413).
\end{solution}

\begin{solution}{exo:b3:computational-chemistry:4}
Acht atomen geven $3N - 6 = 18$ vibraties: 17 reële en één imaginaire. Eén
negatieve \omterm{def:b3:quantum-model-systems:operator}{eigenwaarde} van de hessematrix: een zadelpunt van de eerste
orde, een overgangsstructuur (hier bijvoorbeeld van een rotatie of een
eliminatie).
\end{solution}

\begin{solution}{exo:b3:computational-chemistry:5}
$\zeta = 27/16$, $E = -(27/16)^2 = \qty{-2.8477}{\hartree} = \qty{-77.49}{eV}$.
Fout $-2.8477 + 2.9034 = \qty{0.0557}{\hartree} = \qty{1.52}{eV}$, dus
1.9\,\% van de totale energie --- maar groter dan veel reactie-energieën.
\end{solution}

\begin{solution}{exo:b3:computational-chemistry:6}
$(-1 - E)(-0.5 - E) - (-0.2 - 0.3E)^2 = 0$, dus $0.91E^2 + 1.38E + 0.46 = 0$:
$E = \qty{-1.0217}{\hartree}$ en \qty{-0.4947}{\hartree}. Ja: het bijmengen
van de tweede functie verlaagt de energie onder $H_{11}$, zoals het
\omterm{thm:b3:computational-chemistry:variational}{variatieprincipe} toelaat.
\end{solution}

\begin{solution}{exo:b3:computational-chemistry:7}
$(102/36)^4 = 64$: ongeveer 640 minuten, zo'n 11 uur.
\end{solution}

\begin{solution}{exo:b3:computational-chemistry:8}
$0.5782 \times 27.211 = \qty{15.73}{eV}$, 0.30 eV boven de gemeten waarde.
Bevroren orbitalen: het echte ion relaxeert en ligt lager, dus de waarde
van Koopmans is te hoog. Ontbrekende correlatie: het neutrale molecuul
(twee elektronen) heeft meer \omterm{def:b3:computational-chemistry:correlation}{correlatie-energie} dan het ion (één
elektron), wat de echte waarde groter maakt.
Hier wint de eerste fout.
\end{solution}

\begin{solution}{exo:b3:computational-chemistry:9}
$-0.5960 - 2(-0.4666) = \qty{0.3372}{\hartree} = \qty{9.18}{eV}$ boven twee
atomen. De RHF-functie behoudt 50\,\% \ce{H+ H-}, en het scheiden van een
proton en een hydride kost het verschil tussen de ionisatie-energie van H
en de elektronenaffiniteit van H, dat in de energie meegemiddeld wordt.
\end{solution}

\begin{solution}{exo:b3:computational-chemistry:10}
$E(c_1^2 + 2c_1c_2S + c_2^2) = c_1^2H_{11} + 2c_1c_2H_{12} + c_2^2H_{22}$.
Differentiëren naar $c_1$ met $\partial E/\partial c_1 = 0$ geeft
$E(2c_1 + 2c_2S) = 2c_1H_{11} + 2c_2H_{12}$, dus $(H_{11} - E)c_1 + (H_{12} -
ES)c_2 = 0$; op dezelfde manier $(H_{12} - ES)c_1 + (H_{22} - E)c_2 = 0$.
\end{solution}

\begin{solution}{exo:b3:computational-chemistry:11}
$k \approx [E(1.30) - 2E(1.35) + E(1.40)]/(0.05)^2 = 0.001417/0.0025 =
\qty{0.567}{\hartree\per\bohr\squared} = \qty{882}{N/m}$. Met $\mu = m_H/2 =
\qty{8.37e-28}{kg}$ is $\omega = \sqrt{k/\mu} = \qty{1.03e15}{s^{-1}}$ en
$\tilde\omega = \omega/2\pi c = \qty{5452}{cm^{-1}}$, 24\,\% boven de gemeten
$\tilde\omega_e$: een RHF-kromme in een minimale basis is te steil (geen
correlatie, een te kleine basis); daarom worden berekende frequenties naar
beneden geschaald.
\end{solution}

\begin{solution}{exo:b3:computational-chemistry:12}
Verschillen van enkele \unit{kJ/mol} waarbij een waterstofbrug betrokken is,
vereisen correlatie en dispersie: een hybride functionaal met een
dispersiecorrectie, of beter een gecorreleerde golffunctiemethode voor de
uiteindelijke energieën, met een gepolariseerde drievoudig gesplitste
basis met \omterm{def:b3:computational-chemistry:split-valence}{diffuse functies}. Optimaliseer beide conformeren op
hetzelfde niveau, bevestig de minima met frequenties (alle reëel), tel de
nulpuntsenergieën op, test de basis met één grotere berekening, en
vergelijk met een verwant systeem waarvan de conformatie-energie bekend is.
\end{solution}

\begin{solution}{pb:b3:computational-chemistry:1}
\textbf{1.} Elke \omterm{def:b3:computational-chemistry:basis}{slaterorbitaal} van een minimale basis wordt vervangen door
een vaste contractie van drie gaussfuncties die erop aangepast zijn.
\textbf{2.} De heliumkern trekt haar elektronen sterker aan, dus haar
$1s$-functie is compacter (grotere exponenten). $6.3624/3.4253 = 1.857 = (1.69/1.24)^2$.
\textbf{3.} Twee basisfuncties, twee elektronen, één bezette orbitaal (en
één virtuele).
\textbf{4.} $V_{\mathrm{nn}} = Z_{\mathrm{He}}Z_{\mathrm H}/R = 2/1.4632 =
\qty{1.3669}{\hartree}$.
\textbf{5.} De twee functies overlappen sterk (54\,\%): de atomen zijn
dicht genoeg bij elkaar om een binding te vormen.
\textbf{6.} $h_{11}$ is de energie van een elektron in de functie van He
met beide kernen maar zonder ander elektron; de kern van He (lading 2)
houdt het veel steviger vast.
\textbf{7.} $(h_{11} - \varepsilon)(h_{22} - \varepsilon) - (h_{12} - \varepsilon S)^2 = 0$:
$0.71185\varepsilon^2 + 2.79292\varepsilon + 2.44969 = 0$.
\textbf{8.} $\varepsilon = \qty{-2.600}{\hartree}$ en \qty{-1.324}{\hartree}.
\textbf{9.} $\mathbf h$ verwaarloost de afstoting tussen de twee
elektronen; de core-orbitaal is te gecontraheerd en ligt te laag.
\textbf{10.} $F_{\mu\nu} = h_{\mu\nu} + \sum_{\lambda\sigma}P_{\lambda\sigma}[(\mu\nu|\sigma\lambda) -
\frac12(\mu\lambda|\sigma\nu)]$, met $P_{\lambda\sigma} = 2C_{\lambda1}C_{\sigma1}$.
\textbf{11.} De coulombafstoting van elk elektron door de ladingswolk van
het andere (en, in het algemeen, de uitwisseling).
\textbf{12.} $\mathbf F$ hangt af van $\mathbf P$, dat afhangt van de orbitalen
die uit $\mathbf F$ volgen: de vergelijkingen zijn niet-lineair en worden
opgelost door opeenvolgende benaderingen tot zelfconsistentie bereikt is.
\textbf{13.} $c_1^2 + c_2^2 + 2c_1c_2S = 0.7684 + 0.0410 + 0.1906 = 1.0000$.
\textbf{14.} He: $1.5369 + 0.1906 = 1.727$; H: $0.0820 + 0.1906 = 0.273$
elektron.
\textbf{15.} Ladingen: He $+0.27$, H $+0.73$. De positieve lading zit
grotendeels op waterstof: \ce{HeH+} is een proton gebonden aan een
heliumatoom, \ce{He + H+}, zoals verwacht, want de ionisatie-energie van He
(24.6 eV) is veel groter dan die van H (13.6 eV).
\textbf{16.} $-\varepsilon_1 = \qty{1.633}{\hartree} = \qty{44.4}{eV}$.
\textbf{17.} $E_{\mathrm{el}} = -2.8418 - 1.3669 = \qty{-4.2087}{\hartree}$.
\textbf{18.} Drie cycli bereiken \qty{-2.8418}{\hartree}. Elke cyclus geeft
een determinant waarvan de energie een bovengrens is voor de
geconvergeerde hartree-fockenergie (\omterm{thm:b3:computational-chemistry:variational}{variatieprincipe}), en de iteraties
verbeteren hem.
\textbf{19.} De basis met $\zeta = 2.0925$ geeft de laagste energie en is
dus de betere voor dit molecuul (\omterm{thm:b3:computational-chemistry:variational}{variatieprincipe}): de heliumfunctie
is in \ce{HeH+} compacter dan in het vrije atoom, waarvan de standaardbasis
de $\zeta = 1.69$ overneemt.
\textbf{20.} De lege antibindende $\sigma^*$-orbitaal; haar energie zou bij
benadering min de elektronenaffiniteit van \ce{HeH+} zijn (slecht, in zo'n
kleine basis).
\textbf{21.} Nul (geen elektron); $\qty{-2.8078}{\hartree}$.
\textbf{22.} $-2.8078 - (-2.8418) = \qty{0.0340}{\hartree} = \qty{89}{kJ/mol}$.
\textbf{23.} Ongeveer de helft van de gemeten waarde: een minimale basis
kan de dichtheid van helium niet naar het proton toe polariseren (er zijn
geen $p$-functies), dus de binding is te zwak. (Correcties voor de
\omterm{def:b3:quantum-model-systems:oscillator}{nulpuntsenergie} en voor \qty{298}{K} zijn daarnaast klein.)
\textbf{24.} $-(24.5874 + 54.4178)/27.211 = \qty{-2.9034}{\hartree}$; het
STO-3G-atoom ligt \qty{0.0956}{\hartree} (\qty{2.6}{eV}) te hoog.
\textbf{25.} Geometrieën hangen af van de manier waarop de energie
verandert met $R$; het grootste deel van de fout, geconcentreerd in de
kernschil, is bij elke geometrie vrijwel hetzelfde en valt weg.
\textbf{26.} \textbf{$E(\ce{HeH+}) = \qty{-2.8418}{\hartree}$ op RHF/STO-3G-niveau
(standaardbasis), bij \qty{1.4632}{\bohr}} (en \qty{-2.8607}{\hartree} met de
klassieke $\zeta_{\mathrm{He}} = 2.0925$).
\end{solution}
